\section*{Chapter \ref{ch:qm:monte-carlo} --- Monte Carlo}

\begin{solution}{exo:qm:monte-carlo:1}
Error scales as $n^{-1/2}$: a factor of four needs $16 \times 100\,000 = 1.6$ million paths, or a variance-reduction factor of 16 at 100\,000 paths.
\end{solution}

\begin{solution}{exo:qm:monte-carlo:2}
$1/(1 - \rho^2)$: 5.3 at $\rho = 0.9$ and 50.3 at $\rho = 0.99$. High correlations are needed for large factors.
\end{solution}

\begin{solution}{exo:qm:monte-carlo:3}
$\Var(\frac12(f(Z) + f(-Z))) = \frac12\Var f(Z)(1 + \mathrm{Corr}(f(Z), f(-Z)))$ per pair, i.e.\ per two evaluations; it beats plain sampling when the correlation is negative, which a monotone $f$
guarantees. For an even function $f(Z) = f(-Z)$, the correlation is 1 and the pair is worth one draw at the cost of two.
\end{solution}

\begin{solution}{exo:qm:monte-carlo:4}
By the law of total variance, $\sigma^2 = \sum_sp_s\sigma_s^2 + \sum_sp_s(\mu_s - \mu)^2$. The stratified \omterm{def:qm:estimation:estimator}{estimator} $\sum_sp_s\bar Y_s$ with $n_s = np_s$ has variance $\sum_sp_s^2\sigma_s^2/n_s =
\frac1n\sum_sp_s\sigma_s^2$, which drops the nonnegative between-strata term.
\end{solution}

\begin{solution}{exo:qm:monte-carlo:5}
The \omterm{def:qm:estimation:estimator}{estimator} is $\mathbf 1\{Z > z^*\}e^{-\theta Z + \theta^2/2}$ with $Z \sim \mathcal N(\theta, 1)$; its second moment is
\[\E_f\bigl[\mathbf 1\{Z > z^*\}e^{-\theta Z + \theta^2/2}\bigr] = e^{\theta^2}\bar\Phi(z^* + \theta).\]
The derivative of its logarithm is $2\theta - \varphi(z^* + \theta)/\bar\Phi(z^* + \theta) \approx 2\theta - (z^* + \theta)$ for large arguments (Mills's ratio), which vanishes at $\theta \approx z^*$. At $z^* = 3.67$
the exact minimiser is 3.80, where the simulation's minimum is.
\end{solution}

\begin{solution}{exo:qm:monte-carlo:6}
With $a(x) = \mu x$ and $b(x) = \sigma x$, $bb' = \sigma^2x$: $X_{k+1} = X_k(1 + \mu h + \sigma\Delta W + \frac12\sigma^2(\Delta W^2 - h))$. Expanding the exact step,
$e^{(\mu - \sigma^2/2)h + \sigma\Delta W} = 1 + \sigma\Delta W + (\mu - \frac12\sigma^2)h + \frac12\sigma^2\Delta W^2 + O(h^{3/2})$, the same terms.
\end{solution}

\begin{solution}{exo:qm:monte-carlo:7}
With $2^{14}$ paths: correlation 0.99953, 0.99941 and 0.99425, variance factors 1\,071, 853 and 87 at strikes 80, 100 and 130. Far out of the money the geometric average (always below the
arithmetic one) finishes in the money less often than the arithmetic one, and the two payoffs decouple near the strike: the control is worth less where both are rare.
\end{solution}

\begin{solution}{exo:qm:monte-carlo:8}
Sobol points are not independent, so $s/\sqrt n$ is not the \omterm{def:qm:estimation:estimator}{standard error} of their average: it measures the spread of the payoffs, not of the estimate, and says nothing about the error, which
may be smaller or, with a poor path construction, larger. The error must come from independent scrambles (replicates), and unscrambled points give no error estimate at all.
\end{solution}

\begin{solution}{pb:qm:monte-carlo:1}
\textbf{1.} 6.495 with \omterm{def:qm:estimation:estimator}{standard error} 0.077 from $2^{14}$ paths; reference 6.4788 (\omterm{def:qm:monte-carlo:vr}{control variate} plus scrambled Sobol, 32 replicates of $2^{14}$, \omterm{def:qm:estimation:estimator}{standard error} 0.00007).
\textbf{2.} 16.
\textbf{3.} At equal accuracy the cost is $c\sigma^2/\epsilon^2$, so \omterm{def:qm:estimation:estimator}{estimators} are compared by variance times cost per draw; all the factors here are at equal payoff evaluations.
\textbf{4.} Each path's numbers are a function of (seed, trade, path, block), not of a sequential state, so any split of the paths reproduces them.
\textbf{5.} It maps one uniform coordinate to one normal monotonically, so the structure of the point set survives; Box--Muller mixes two coordinates through a cosine.
\textbf{6.} 1.8: the two payoffs of a pair have correlation $-0.43$; the payoff is monotone in the path but convex, and its positive part discards the down paths.
\textbf{7.} The geometric-average call, 6.1793 in closed form; correlation 0.99941; factor 853.
\textbf{8.} $\beta^* = \Cov(Y, C)/\Var C$; estimated from the same paths as 1.04 (the bias this introduces is of order $1/n$).
\textbf{9.} Where the payoff lives in a small region of the driver's space: at $\theta = 3.8$ the relative \omterm{def:qm:estimation:estimator}{standard error} falls from 20\% to 0.64\%, a factor of about 1\,000 in variance.
\textbf{10.} The likelihood ratio $e^{-\theta Z + \theta^2/2}$ becomes very variable: the second moment $e^{\theta^2}\bar\Phi(z^* + \theta)$ grows again beyond the optimum.
\textbf{11.} About 34 with the increments in time order and about 2\,500 with the bridge.
\textbf{12.} Low-discrepancy points are most uniform in their first coordinates; the bridge puts $W_T$ and the coarse shape of the path there, the directions that carry most of the payoff's variance.
\textbf{13.} About 37\,000 (\omterm{def:qm:estimation:estimator}{standard error} 190 times smaller); fitted slopes between $-0.7$ and $-0.8$ against $-0.5$ for Monte Carlo.
\textbf{14.} From the spread of independent scrambles: $R$ replicates give $R$ independent unbiased estimates, whose sample standard deviation divided by $\sqrt R$ is the \omterm{def:qm:estimation:estimator}{standard error}.
\textbf{15.} Strong: the pathwise error, slopes 0.49 (Euler) and 0.89 (Milstein); weak: the error in law, slope 1.00 for the mean.
\textbf{16.} Cost $3.2 \times 10^7$ time steps against $3.0 \times 10^9$, a factor of 94: most samples are drawn on coarse, cheap levels, and the fine corrections have small variance because Milstein's
strong order 1 makes $V_l$ fall by four per level.
\textbf{17.} No: it relies on a smooth payoff with a closed-form control and a low effective dimension. Barriers, digitals and early exercise break the smoothness, and most trades have no control
as good; each trade class must be measured.
\textbf{18.} Seed, trade and path counters, generator and version, scrambling seeds, the path construction, and the point-set size.
\textbf{19.} Named result: \emph{sixteen times fewer machines}: on the representative trade, the \omterm{def:qm:monte-carlo:vr}{control variate} alone reduces the variance by 853 and with scrambled Sobol points and the \omterm{def:qm:brownian-motion:bridge}{Brownian bridge} by about
37\,000, against the 16 required; the target is met at today's number of paths, with the generation overhead to be measured in the production engine.
\textbf{20.} Into the constant in front of $n^{-1/2}$: what each path is worth, and what it costs.
\end{solution}

\begin{solution}{iq:qm:monte-carlo:1}
The \omterm{def:qm:estimation:estimator}{standard error} is $\sigma/\sqrt n$ whatever the dimension: the dimension enters only through $\sigma$ and the cost of a path, which is why Monte Carlo wins in high dimension and loses to grids in one or two.
\iqlookfor{$n^{-1/2}$, dimension-free, the constant.}
\end{solution}

\begin{solution}{iq:qm:monte-carlo:2}
\omterm{def:qm:monte-carlo:vr}{Antithetic variates} (monotone payoffs, cheap, small gains); \omterm{def:qm:monte-carlo:vr}{control variates} (a correlated quantity with known mean: the geometric Asian, the underlying, a Black--Scholes price); \omterm{def:qm:monte-carlo:is}{importance sampling}
(rare events: shift the drift towards the region that matters); stratification or \omterm{def:qm:monte-carlo:qmc}{quasi-Monte Carlo} (smooth payoffs in a few important directions).
\iqlookfor{When each works, not just names.}
\end{solution}

\begin{solution}{iq:qm:monte-carlo:3}
Use a \omterm{def:qm:monte-carlo:prng}{counter-based generator} (Philox, Threefry) keyed by the seed, with the counter built from trade, path and block indices, so that numbers are a function of their coordinates, not of the order
of generation; reduce partial sums in a fixed order.
\iqlookfor{Counters over sequential streams, and a fixed reduction order.}
\end{solution}

\begin{solution}{iq:qm:monte-carlo:4}
Its uniformity is good only in the first coordinates, and in time order the path's variance is spread evenly over 250 of them. Fix it with a \omterm{def:qm:brownian-motion:bridge}{Brownian bridge} or principal-component construction that
puts the important directions first, and scramble to get error estimates.
\iqlookfor{Effective dimension and path construction.}
\end{solution}

\begin{solution}{iq:qm:monte-carlo:5}
\omterm{def:qm:monte-carlo:is}{Importance sampling}: tilt the distribution (shift the mean of the drivers, or exponentially tilt the loss) so that the loss beyond the threshold is common, and weight by the likelihood ratio; choose the
tilt near the most likely point of the rare region; check the weights' spread. Splitting or conditional Monte Carlo are alternatives.
\iqlookfor{A \omterm{def:qm:probability-at-speed:change}{change of measure} with the likelihood ratio, and a sensible choice of tilt.}
\end{solution}

\begin{solution}{iq:qm:monte-carlo:6}
Write the finest-level expectation as a telescoping sum of corrections between successive step sizes, simulate the fine and coarse paths of each correction from the same Brownian path, and allocate
samples in proportion to $\sqrt{V_l/C_l}$. The correction's variance $V_l$ is governed by the strong error, so a higher strong order makes the fine levels need very few samples; with $V_l$ falling
faster than the cost rises, the total cost is $O(\epsilon^{-2})$, as for an \omterm{def:qm:estimation:estimator}{unbiased estimator}.
\iqlookfor{Telescoping sum, coupling, and the variance of corrections from the strong order.}
\end{solution}
